% Vorlage f�r WABs der Provadis Hochschule
% basiert auf "Universit�t Ulm Praktikumsbericht Vorlage"
% von Max Sch.
% CC BY 4.0
%
% main.tex
%
% Das ist das zentrale Dokument.


%pers�nliche Infos
% Vorlage für WABs der Provadis Hochschule
% basiert auf "Universität Ulm Praktikumsbericht Vorlage"
% von Max Sch.
% CC BY 4.0
%
% Berichtsinfos.tex
%
% Hier werden individuelle Inforamtionen festgelegt.

% Art der Arbeit (genau eine Zeile aktiv lassen)
% Pflichtangabe Deckblatt, siehe Anforderungen "Formatierung 1a"
\newcommand{\artderarbeit}{Wissenschaftlich angeleitete Berufspraxis}
%\newcommand{\artderarbeit}{Wissenschaftlicher Kurzbericht}
%\newcommand{\artderarbeit}{Bachelorarbeit}
%\newcommand{\artderarbeit}{Masterarbeit}

% Nur bei Praxisberichten (WAB / Kurzbericht): das Fach. Bei BT/MT leer lassen: {}
\newcommand{\fach}{Name des Fachs / der Veranstaltung}

% Nur bei BT/MT: angestrebter akademischer Grad. Bei Praxisberichten leer lassen: {}
\newcommand{\abschluss}{}
%\newcommand{\abschluss}{Bachelor of Science (B.Sc.)}
%\newcommand{\abschluss}{Master of Science (M.Sc.)}

%persönliche Daten
\newcommand{\fullname}{First Last}
\newcommand{\email}{first.last@stud-provadis-hochschule.de}
\newcommand{\matnr}{XXXX}
\newcommand{\studiengang}{Business Information Management}
\newcommand{\studiengruppe}{BBIM26}
% Telefonnummer ist optional. Zum Weglassen leer lassen: {}
\newcommand{\telMobil}{+49 XXX XXXXXXX}
\newcommand{\fakultaet}{Informatik und Wirtschaftsinformatik}

%Daten zum Unternehmen
\newcommand{\titel}{Textvorlage für wissenschaftliche Arbeiten\\
Titel und Untertitel der Arbeit}
\newcommand{\unternehmen}{Musterfirma}

\newcommand{\abgabedatum}{XX.XX.XXXX}
\newcommand{\ort}{Frankfurt am Main}

% Optional: Bild der digitalen Unterschrift für die Selbstständigkeitserklärung,
% z.B. {Bilder/Unterschrift.png}. Bei Upload im HVS optional -> leer lassen: {}
\newcommand{\unterschrift}{}

% Daten zu Betreuern
% Referent:in = betreuende Dozent:in
\newcommand{\refFirst}{Prof. Dr. First Last}
% Korreferent:in nur bei Bachelor- und Masterarbeiten, sonst leer lassen: {}
\newcommand{\refSecond}{Prof. Dr. First Last}

% Zitierweise: Kurzzitierweise nach Harvard in Fußnoten (typisch in der Wirtschaftsinformatik,
% "siehe Weber (2020): S. 8.") oder numerisch (typisch in der Kerninformatik, IEEE-ähnlich,
% "[1, S. 8]").
\newif\ifharvard
\harvardtrue
%\harvardfalse

% Wurde KI verwendet? Steuert den Inhalt des (verpflichtenden) KI-Verzeichnisses.
\newif\ifKIverwendet
\KIverwendettrue
%\KIverwendetfalse


%Deklariert den Namen des Authors, sowie den Namen der generierten pdf Datei
\title{\artderarbeit\ -- \unternehmen}
\author{\fullname}
\date{\abgabedatum}


\documentclass[
	a4paper,
	ngerman,
	12pt,	% 12pt f�r Times, 11pt f�r Arial und 10pt f�r Verdana
	bibliography=totoc,
	listof=totoc,
	numbers=noendperiod,
	openany,
	twoside,
	%oneside, % Anforderung: links 3 cm, rechts 2 cm auf JEDER Seite (twoside spiegelt die R�nder)
	headinclude=false,
	titlepage=firstiscover,
	abstract=true,
	headsepline=true,
	footsepline=true,
	headings=normal
	]{scrreprt}
%Startet das Dokument und regelt gleichzeitig Formatierungen wie Schriftgr��e, Papierformat, Zahlenart, Sprache usw.
% 12pt - die Standardschriftgr��e
% openany - Kapitel d�rfen auch auf der "R�ckseite" starten, Leerseiten werden vermieden
% parskip=full - Leerzeilen zwischen Abs�tzen statt Einr�ckung
% twoside - R�ckseiten werden mit bedruckt, Seitenzahlen und Binder�nde werden entsprechen angepasst
% Alternative: oneside - einseitiges Layout, damit der linke Rand auf jeder Seite 3 cm betr�gt

% \RedeclareSectionCommand[
%   beforeskip=-32pt
% ]{chapter}
% Optionaler Befehl zum editieren von dem margin top von Kapiteln

%allgemeine Einstellungen
% Vorlage für WABs der Provadis Hochschule
% basiert auf "Universität Ulm Praktikumsbericht Vorlage"
% von Max Sch.
% CC BY 4.0
%
% Einstellungen.tex
%
% Hier werden allgemeine Einstellungen festgelegt.

%Packages sind Erweiterungen des eigentlichen Programms und erleichtern Abläufe, Darstellungen, Formatierungen usw.
%Im Folgenden sind die Packages installiert, welche ich zum erstellen einer Bachelorarbeit einmal benötigt wurden.
%Es kann beliebig erweitert/reduziert werden.

%%%%%%%%%%%%%%%%%%%%%%%%%%%%%%%%%%%%%%%%%%%%%%%%%%%%%%%%%
%%%%%%%%%%%%%%%%%%%%%%%%%%%%%%%%%%%%%%%%%%%%%%%%%%%%%%%%%

% T1-Schriftkodierung: nötig, damit Wörter mit Umlauten korrekt getrennt werden
% und der Text im PDF durchsuch-/kopierbar ist (u.a. für die Plagiatsprüfung)
\usepackage[T1]{fontenc}
% Times New Roman Alternative, pdflatex-kompatibel
\usepackage{newtxtext}

% Paket für Times Schrift
%\usepackage{mathptmx}

\usepackage[ngerman]{babel}
% Paket für Deutsche Sprache (Übersetzungen von Chapter zu Kapitel,
% richtige Umlaute, richtige % Silbentrennung)
% siehe auch http://de.wikipedia.org/wiki/Babel-System

\usepackage{titling}
%Hilft im Dokument auf den deklarierten Autor und Titel zuzugreifen

\usepackage{csquotes}
% Zitierweise wird in Berichtinfos.tex über \harvardtrue / \harvardfalse gewählt.
% Anforderung: einheitlicher Stil, Seitenangabe obligatorisch; Online-PDFs mit URL oder DOI
% (daher doi=true).
\ifharvard
    % Kurzzitierweise nach Harvard in Fußnoten, wie im Anforderungsdokument empfohlen:
    %   Fußnote:              siehe Weber (2020): S. 8.
    %   Literaturverzeichnis: Weber, Max (2020): Titel. Ort: Verlag.
    %   Sortierung:           alphabetisch nach Verfasser, dann Jahr, dann Titel (nyt)
    % Verwendung im Text:     \footcite[Vgl.][8]{schluessel}   bzw. \footcite[8\psq]{schluessel} (=S. 8 f.)
    \usepackage[style=authoryear,backend=biber,sorting=nyt,doi=true,isbn=false,
        maxcitenames=2,maxbibnames=3,dashed=false,uniquename=false,uniquelist=false]{biblatex}
    \DeclareDelimFormat{nameyeardelim}{\addspace}
    \DeclareDelimFormat{postnotedelim}{\addcolon\space}
    \DeclareDelimFormat[bib]{nameyeardelim}{\addspace}
    \DeclareDelimFormat[bib]{labelnamepunct}{\addcolon\space}
    \renewbibmacro*{cite:labeldate+extradate}{%
        \iffieldundef{labelyear}{}%
            {\printtext[bibhyperref]{\printtext[parens]{\printlabeldateextra}}}}
\else
    % Numerische Zitierweise (IEEE-ähnlich), Seitenangabe trotzdem Pflicht: \cite[8]{schluessel}
    \usepackage[style=numeric-comp,backend=biber,doi=true,isbn=false,maxnames=3,sorting=none]{biblatex}
\fi
% Internetquellen: Webadresse in eckigen Klammern plus Datum des letzten Zugriffs
\DeclareFieldFormat{url}{\mkbibbrackets{\url{#1}}}
\DefineBibliographyStrings{ngerman}{urlseen = {Zugriff am}}
\addbibresource{Referenzen.bib}
%Regelt die Erstellung des Literaturverzeichnisses sowie die Zitation im gesamten Bericht
%Greift gleichzeitig auf die Referenzen.bib genannte BibTex datei zu, in welcher die Quellen aufgeführt sind.

\usepackage{url}
\urlstyle{same}
%erzeugt schönere URLs (vorallem in den Quellen wichtig), in der Textschrift statt Schreibmaschinenschrift

\usepackage{float}
%Hilft beim Anordnen von Bildern, Grafiken und Tabellen

%\usepackage{abstract}
%ermöglicht das einfügen eines Abstract (Nicht relevant fürs Praktikum)

\usepackage{booktabs}
%ermöglicht das erstellen von schöneren Tabellen

\usepackage{array}
\usepackage{longtable}
%Tabellen, die über mehrere Seiten umbrechen dürfen (z.B. KI-Verzeichnis)

\usepackage{subfigure}
%lässt den Benutzer mehrere Grafiken/Bilder in einer Abbildung darstellen

\usepackage{setspace}
% Paket um den Zeilenabstand zu ändern
\onehalfspacing
% Zeilenabstand auf 1,5-fach setzen
% Zeilenabstand 1,5 Zeilen für den Fließtext (Anforderung).
% Hinweis: \onehalfspacing = 1,5 x Schriftgröße (18 pt bei 12 pt). Wer exakt Words
% "1,5 Zeilen" bei Times New Roman 12 (ca. 20,7 pt) nachbilden will: \setstretch{1.43}

\usepackage{pdfpages}
%ermöglicht das einbinden von eigenständigen PDF Dokumenten

\usepackage{siunitx}
\sisetup{locale = DE}
%vereinfacht den benötigten Syntax zum Benutzen von SI Einheiten und Mathematischen Gleichungen mit dem deutschen Standart

\usepackage{tikz}
%Freies Anordnen von Bildern und Text in Relation zueinander, sowie Möglichkeiten Bilder zu bearbeiten
%Verfügt auserdem über die Möglichkeit Grafiken und Diagramme in Latex zu erstellen

\usepackage[format=plain,labelfont=bf]{caption}
\captionsetup[table]{position=above}
%verschönert die Abbildungs- und Tabellenbeschriftungen. Für Tabellen werden die Beschriftungen über der Tabelle angezeigt

\usepackage{icomma}
%Für den korrekten Abstand der Kommas in Dezimalzahlen und in Gleichungen

\usepackage{amsmath}
%elementares Paket für Gleichungen. Verschönert diese und erleichtert die Benutzung von Gleichungen.

\usepackage{graphicx}
%mehr Optionen beim Einbinden und anordnen von Bildern und Grafiken

%%% HJP 31/3/25
\newlength\dist
\setlength\dist{-\topskip}
\let\chapback\chapterheadstartvskip
\def\setchapspace#1{\setlength\dist{#1\dist}}
\def\lesschapspace{\renewcommand*\chapterheadstartvskip{\vspace*{\dist}}}
\def\resetchapspace{\let\chapterheadstartvskip\chapback}
%%%

\usepackage[left=3cm,right=2cm,top=3cm,bottom=2cm,footskip=1cm]{geometry}
%Regelt die Seitenabstände gemäß Anforderung: links und oben je 3 cm, rechts und unten je 2 cm
%(setzt einseitiges Layout voraus, siehe Option oneside in main.tex)

\usepackage{enumitem}

\usepackage{etoolbox}
%Ermöglicht das Entfernen des Seitenumbruchs vor \chapter (siehe Abschnitt unten)

\usepackage[hidelinks]{hyperref}
% Generische Links

\usepackage[acronym, toc]{glossaries}

\renewcommand*{\glspostdescription}{}
%Entfernt die Nummerierung im Abkürzungsverzeichnis

\usepackage{listings}
\definecolor{codegreen}{rgb}{0,0.6,0}
\definecolor{codegray}{rgb}{0.5,0.5,0.5}
\definecolor{codepurple}{rgb}{0.58,0,0.82}
\definecolor{backcolour}{rgb}{0.95,0.95,0.92}
\lstdefinestyle{myStyle}{
    backgroundcolor=\color{backcolour},
    commentstyle=\color{codegreen},
    keywordstyle=\color{magenta},
    numberstyle=\tiny\color{codegray},
    stringstyle=\color{codepurple},
    basicstyle=\ttfamily\footnotesize,
    breakatwhitespace=false,
    breaklines=true,
    keepspaces=true,
    numbers=left,
    numbersep=5pt,
    showspaces=false,
    showstringspaces=false,
    showtabs=false,
    tabsize=2,
}

\usepackage{todonotes}
% Kann mit disable komplett ausgeschaltet werden
%\usepackage[disable]{todonotes}
\setuptodonotes{fancyline}
\todototoc

\flushbottom

\iffalse % don’t use titlesec.sty with KOMA script classes
\usepackage{titlesec}
%% \titleformat{⟨command⟩}[⟨shape⟩]{⟨format⟩}{⟨label⟩}{⟨sep⟩}{⟨before⟩}[⟨after⟩]
\titleformat{name=\chapter}[display]
{\raggedleft\chaptertitlename\ {\textcolor{gray}{\fontsize{60}    {65}\selectfont\thechapter}}}
{20pt}{}[]
\fi

%%%%%%%%%%%%%%%%%%%%%%%%%%%%%%%%%%%%% ToC DEPTH LEVEL
\setcounter{secnumdepth}{3} % number subsubsection
\setcounter{tocdepth}{3} % list subsubsection


\addtokomafont{chapterprefix}{\raggedleft}

%%%%%%%%%%%%%%%%%%%%%%%%%%%%%%%%%%%%% KAPITEL: BEGINN AUF NEUER SEITE?
% KOMA-Script beginnt jedes Kapitel auf einer neuen Seite (\clearpage, bei "openright" zusätzlich \cleardoublepage).
%
% Schalter:
%   \kapitelneuseitefalse  Die Kapitel laufen fort: mehrere Kapitel koennen
%                          auf derselben Seite stehen. Der Seitenumbruch vor
%                          \chapter wird dazu entfernt - Nummerierung, Format
%                          und Inhaltsverzeichnis der Kapitel bleiben
%                          unveraendert.
%   \kapitelneuseitetrue   Jedes Kapitel beginnt auf einer neuen Seite
%
% Die Kapitel werden in main.tex ueber \kapitel{...} eingebunden: im Standard
% mit \include (eigene .aux-Datei je Kapitel, Seitenumbruch am Dateiende), im
% Fliesstext mit \input (kein erzwungener Seitenumbruch am Dateiende).
\newif\ifkapitelneuseite
\kapitelneuseitetrue
% \kapitelneuseitefalse

\ifkapitelneuseite
  % Standard: Jedes Kapitel beginnt auf einer neuen Seite - nichts zu tun.
\else
  % Fliesstext: Der Seitenumbruch vor \chapter wird entfernt.
  \makeatletter

  \patchcmd{\scr@startchapter}
    {\if@openright\cleardoublepage\else\clearpage\fi}
    {}
    {\relax}
    {%
      \PackageError{WAB-Vorlage}{%
        Der Seitenumbruch vor \string\chapter konnte nicht entfernt werden%
      }{%
        Die interne Definition von \string\scr@startchapter (KOMA-Script) hat%
        sich geaendert. Bitte die Vorlage entsprechend anpassen.%
      }%
    }

  \renewcommand*{\scr@chapter@beforeskip}{1.5\baselineskip}%

  \providecommand*{\scr@chapter@afterindent}[3]{\@afterindentfalse}%

  \renewcommand*{\chapterheadstartvskip}{\vspace{\@tempskipa}}%
  \makeatother
\fi

% Kapitel-Dateien einbinden: \kapitel{Kapitel/1Einleitung}
%   Fliesstext: \input (kein erzwungener Seitenumbruch am Dateiende, damit
%               mehrere Kapitel auf derselben Seite stehen koennen)
%   Standard:   \include (eigene .aux-Datei je Kapitel, Seitenumbruch am
%               Dateiende, daher beginnt jedes Kapitel auf einer neuen Seite)
\ifkapitelneuseite
  \newcommand*{\kapitel}[1]{\include{#1}}
\else
  \newcommand*{\kapitel}[1]{\input{#1}}
\fi

%%%%%%%%%%%%%%%%%%%%%%%%%%%%%%%%%%%%% PREFACE AND BODY STYLING
\def\preface{
    \pagenumbering{roman}
}

\def\body{
    %\cleardoublepage
    \ifkapitelneuseite
      \lesschapspace
    \fi
    \pagenumbering{arabic}
    \pagestyle{headings}
}

\def\abstract{
    \begin{center}{
        \large\bfseries Kurzzusammenfassung}
    \end{center}
    \normalsize
    \normalfont
}
\def\endabstract{
  \par
}

\newenvironment{acknowledgements}{
   \cleardoublepage
    \begin{center}{
        \large \bf Danksagung}
    \end{center}
    \normalsize
    }{\cleardoublepage}
\def\endacknowledgements{
  \par
}

%%%%%%%%%%%%%%%%%%%%%%%%%%%%%%%%%%%%%%%%%%%%%%%%%%%%%%%%%
%%%%%%%%%%%%%%%%%%%%%%%%%%%%%%%%%%%%%%%%%%%%%%%%%%%%%%%%%

%%%%%%%%%%%%%%%%%%%%%%%%%%%%%%%%%%%%% ZITATE AUS KI-QUELLEN
% Anforderung: Wörtliche KI-Zitate *in Asterisken* statt Anführungszeichen,
% Fußnote verweist auf den Eintrag im KI-Verzeichnis (z.B. "ChatGPT 1.").
% Paraphrasen / KI als Ausgangspunkt: Fußnote "Vgl. Elicit 1."
% Übernahme von Strukturen: Fußnote an der Kapitelüberschrift.
%   \kizitat{Text der KI}{ChatGPT 1}
%   \kivgl{Elicit 1}
\newcommand{\kizitat}[2]{*#1*\footnote{#2.}}
\newcommand{\kivgl}[1]{\footnote{Vgl. #1.}}


\input{Glossar}

\input{Abkuerzungen}

\begin{document}

% Deckblatt erh�lt eine eigene (unsichtbare) Nummerierung, damit es nicht mit
% Seite 1 des Textteils kollidiert (doppelte PDF-Sprungmarke "page.1")
\pagenumbering{Alph}
% Vorlage für WABs der Provadis Hochschule
% basiert auf "Universität Ulm Praktikumsbericht Vorlage"
% von Max Sch.
% CC BY 4.0
%
% Titelseite.tex
%
% Hier wird die Titelseite erzeugt.
% Pflichtangaben laut Anforderungen (Formatierung 1a-f): Art der Arbeit, Titel,
% Fach (nur Praxisberichte), Name, Matrikelnummer, Studiengang mit Studiengruppe,
% E-Mail, Telefon (optional), Abgabedatum, Referent:in (+ Korreferent:in bei BT/MT).
% Alle Werte werden in Berichtinfos.tex gesetzt.
\begin{titlepage}

    % Logo Provadis Hochschule und ggf. Logo Arbeitgeber
    \includegraphics[height=2.06cm]{Bilder/Deckblatt/provadis-hochschule.pdf}
    \hfill
    % optionales Unternehmenslogo
    %\includegraphics[height=2.06cm]{Bilder/Deckblatt/Infraserv_Höchst_logo.pdf}

    \vspace*{1cm}

    \begin{singlespace}
        \begin{center}
            \large
            \artderarbeit

            \vspace*{2cm}

            \Large

            \textbf{\titel}

            \vspace*{2.5cm}

            \normalsize
            Zur Veranstaltung / zur Erlangung des akademischen Grades\\
            ’Bachelor of Science’ B.Sc. / ’Master of Science’ M.Sc.\\
            im Studiengang ’XXX’
            \ifdefempty{\fach}{}{%
                im Fach \enquote{\fach}\\
            }%
            \ifdefempty{\abschluss}{}{%
                zur Erlangung des akademischen Grades\\
                \abschluss\\
            }

            \vspace*{1.5cm}

            vorgelegt dem Fachbereich \fakultaet\ der\\
            Provadis School of International Management and Technology\\
            von

            \vspace*{1cm}

            \fullname \\
            Matrikelnummer: \matnr \\
            Studiengang: \studiengang, Studiengruppe: \studiengruppe \\
            \email
            \ifdefempty{\telMobil}{}{\\ \telMobil}
        \end{center}
    \end{singlespace}

    \normalsize
    \vfill % variabler vertikaler Abstand
    \ifdefempty{\refSecond}%
        {\def\korefzeile{}}%
        {\def\korefzeile{Korreferent:in: & \refSecond \\[0.25cm]}}%
    \begin{tabular}{@{}ll}
        Referent:in: & \refFirst \\[0.25cm]
        \korefzeile
        Abgabedatum: & \abgabedatum \\
    \end{tabular}

\end{titlepage}

% Ende der Datei


\preface
%Ab hier Seitenzahlen mit R�mischen Ziffern

\listoftodos

% Sperrvermerk NUR einkommentieren, wenn zwingend sch�tzensnotwendige,
% vertrauliche Informationen enthalten sind (siehe Kapitel/Sperrvermerk.tex)
%% Vorlage für WABs der Provadis Hochschule
% basiert auf "Universität Ulm Praktikumsbericht Vorlage"
% von Max Sch.
% CC BY 4.0
%
% Sperrvermerk.tex
%
% Wird in main.tex standardmäßig NICHT eingebunden - nur bei zwingend
% schützensnotwendigen Informationen einkommentieren.
\addchap{Sperrvermerk}
\todo[inline]{
Ein Sperrvermerk ist nur zulässig, wenn zwingend \emph{schützensnotwendige}, vertrauliche Informationen in der Arbeit enthalten sind. Dies ist normalerweise nur dann anzunehmen, wenn betriebsinterne Informationen des Arbeitsgebers \emph{ohne Generalisierung, Abstrahierung oder Pseudonymisierung wiedergegeben werden}.\\
Die Überprüfung der Quellen durch die Gutachter muss trotz Sperrvermerk möglich sein. Die mit der Begutachtung und Prüfung von der Hochschule beauftragten Personen (direkt und indirekt) müssen zur Nutzung der Inhalte der Prüfungsleistung berechtigt sein.\\
Eine Nutzung der (digitalen) Dokumente für Zwecke der Plagiatsüberprüfung muss erlaubt sein.\\
Sperrvermerke, die hierzu in Widerspruch stehen, sind wahlweise nichtig oder führen dazu, dass die Arbeit als nicht abgegeben angesehen wird / die Prüfungs\-leistung als nicht erbracht bewertet wird.\\
In begründeten Fällen (siehe oben) kann ein Sperrvermerk wie folgt aussehen:}

\begin{verse}
    Diese Arbeit wurde unter Verwendung von vertraulichen Informationen der Firma \enquote{\unternehmen} angefertigt. Diese Arbeit darf Dritten ohne ausdrückliche Zustimmung der Firma \enquote{\unternehmen} nicht zugänglich gemacht werden.

    Alle Nutzungen für die Begutachtung durch die Hochschule und von ihr beauftragte Dienstleister sind explizit erlaubt. Zu der Begutachtung gehört auch die Prüfung auf Plagiate und die Verwendung künstlicher Intelligenz unter Nutzung entsprechender IT-Plattformen.
\end{verse}


% auskommentieren, um den Abstract/Kurzzusammenfassung wegzulassen
% Vorlage für WABs der Provadis Hochschule
% basiert auf "Universität Ulm Praktikumsbericht Vorlage"
% von Max Sch.
% CC BY 4.0
%
% Abstract.tex

% Im Englischen wird dieser Abschnitt of als ``Abstract'' berzeichnet.
\addcontentsline{toc}{chapter}{Kurzzusammenfassung}
\begin{abstract}
Eine optionale Kurzzusammenfassung.

Bei englischsprachigen Arbeiten sollte auch eine Kurzzusammenfassung in deutscher Sprache angefügt werden.
\end{abstract}


% Reihenfolge der Verzeichnisse gem�� Anforderungen (Formatierung 2a):
% Sperrvermerk, Inhalt, Abbildungen, Glossar, Tabellen, Abk�rzungen, Symbole

\tableofcontents
% erstellt das vollautomatische Inhaltsverzeichnis

\listoffigures
% erstellt das Abbildungsverzeichnis (sofern notwendig)

\listoftables
% erstellt das Tabellenverzeichnis (sofern notwendig)

\glsaddall[types=main]
% Befehl, um alle Glossar Eintr�ge ohne sie im body aufzurufen, auszugeben
\printglossary[title=Glossar, type=\glsdefaulttype]
% erstellt das Glossar (sofern notwendig)

\printglossary[type=\acronymtype, title=Ab�rzungsverzeichnis, nonumberlist]
% erstellt das Abk�rzungsverzeichnis (sofern notwendig)

% Symbolverzeichnis (sofern notwendig), z.B.:
%\addchap{Symbolverzeichnis}
%\begin{tabular}{@{}ll}
%    $\alpha$ & Signifikanzniveau \\
%\end{tabular}

\newpage

\body
%ab hier Seitenzahlen in Arabischen Zahlen

\include{Kapitel/01_Einleitung}
% Vorlage für WABs der Provadis Hochschule
% basiert auf "Universität Ulm Praktikumsbericht Vorlage"
% von Max Sch.
% CC BY 4.0
%
% 02_Stand.tex

\chapter{Stand der Forschung}

In diesem Kapitel wird besonders intensiv mit der Literatur gearbeitet. Entsprechend werden hier sehr viele Zitate vorkommen. Ich werde die Quellenhinweise sofort einbauen, damit ich sie nicht vergesse.

Dabei beachte ich die Vorgaben der Provadis School of International Management and Technology, die in den „Anforderungen an wissenschaftliche Hausarbeiten“ beschrieben sind.

Die Literatur wird im Literaturverzeichnis am Ende der Arbeit aufgeführt. Es enthält alle verwendeten Quellen -- und nur diese --, auch Internetquellen (mit Webadresse in eckigen Klammern und Datum des letzten Zugriffs).

Literaturquellen werden in der Datei Referenzen.bib im BibTeX-Format eingetragen und zitiert. Die Seitenangabe ist obligatorisch.
\ifharvard
    Ein wörtliches Zitat steht in Anführungszeichen: \enquote{Kaffee ist die einzige Substanz, die uns schneller arbeiten lässt}\footcite[12]{lustigeCitation}. Eine Paraphrase steht ohne Anführungszeichen, die Fußnote wird mit \enquote{Vgl.} eingeleitet\footcite[Vgl.][8\psq]{mustermann2020}. Auch Internetquellen werden so belegt\footcite[Vgl.][]{provadisWeb}.
\else
    Ein wörtliches Zitat steht in Anführungszeichen: \enquote{Kaffee ist die einzige Substanz, die uns schneller arbeiten lässt}~\cite[12]{lustigeCitation}. Paraphrasen stehen ohne Anführungszeichen~\cite[8\psq]{mustermann2020}.
\fi

Wörtlich übernommene Stellen aus KI-Quellen stehen in Asterisken: \kizitat{Ein gutes Recruiting beginnt mit einer klaren Stellenbeschreibung.}{ChatGPT 1} Wurde KI-Output nur als Ausgangspunkt verwendet, lautet die Fußnote \enquote{Vgl. Elicit 1.}\kivgl{Elicit 1}

Gut gepflegte BibTeX-Einträge finden sich etwa in der DBLP\footnote{\url{https://dblp.uni-trier.de}} und in den Online-Bibliotheken großer Wissenschaftsverlage, wie der ACM Digital Library\footnote{\url{https://dl.acm.org}}, IEEE Xplore\footnote{\url{https://ieeexplore.ieee.org}}, Springer Link\footnote{\url{https://link.springer.com}}, Elsevier ScienceDirect\footnote{\url{https://www.elsevier.com/de-de/products/sciencedirect}}, uvw.
\todo[inline]{Quellenverweise sollten immer \emph{direkt} hinter Nennung der Quelle bzw. der zu unterstützenden Aussage stehen, und nicht am Ende des Satzes oder gar des Absatzes.}

Weitere Verzeichnisse sind alphabetisch zu sortieren und besonders sorgfältig anzufertigen.

\include{Kapitel/03_Methodik.tex}
\include{Kapitel/04_Untersuchung.tex}
\include{Kapitel/05_Ergebnisse.tex}
\include{Kapitel/06_Fazit.tex}

\pagenumbering{Roman}
% Seiten nach dem Textteil mit gro�en r�mischen Zahlen (I, II, III ...)

% Reihenfolge gem�� Anforderungen (Formatierung 2b):
% Literatur, KI-Verzeichnis, Rechtsquellen, Anhang, Selbstst�ndigkeitserkl�rung

\printbibliography[title=Literaturverzeichnis]
% erstellt das Literaturverzeichnis (verpflichtend)

% Vorlage für Praktikumsberichte
%
% A1_KI.tex
%
% Das KI-Verzeichnis ist verpflichtend. Ob KI verwendet wurde, wird in
% Berichtinfos.tex über \KIverwendettrue / \KIverwendetfalse eingestellt.

\addchap{KI-Verzeichnis}

\todo[inline]{
Jede Nutzung von KI ist zu dokumentieren. Dazu wird hinter dem Literaturverzeichnis ein separates KI-Verzeichnis eingefügt, das alle KI-generierten Inhalte, die eingesetzten Systeme, die verwendeten \emph{Prompts} sowie die weitere Verwendung des Outputs der KI transparent macht.\\
Bei der mehrfachen Verwendung eines Systems werden die Einträge durchnummeriert.\\
Die Reihenfolge der Einträge entspricht der Reihenfolge der Verwendung im Text.\\
In der Spalte \enquote{Prompt} ist der Ausgangsprompt aufzunehmen.\\
Das KI-Verzeichnis ist \textbf{verpflichtend}. Sofern keine KI verwendet wurde, enthält das Verzeichnis nur den Eintrag \enquote{Es wurde keine KI verwendet.} (\texttt{\textbackslash KIverwendetfalse} in Berichtinfos.tex).}

ifKIverwendet
    % longtable statt Gleitumgebung: bleibt direkt unter der Überschrift und
    % darf über mehrere Seiten umbrechen
    {\renewcommand{\arraystretch}{1.5}
    \begin{longtable}{@{}>{\raggedright\arraybackslash}p{0.15\textwidth} >{\raggedright\arraybackslash}p{0.5\textwidth} >{\raggedright\arraybackslash}p{0.27\textwidth}@{}}
        \toprule
        \textbf{System} & \textbf{Prompt} & \textbf{Verwendung (ggf. Beschreibung)}\\
        \midrule
        \endhead
        \bottomrule
        \endfoot
        ChatGPT 1 & What criteria should I use to select a leader? & weiterentwickelt \\
        ChatGPT 2 & Schreibe einen Text, in dem die folgenden Themen behandelt werden: Personalmarketing und seine Bedeutung für ein Unternehmen – der Zusammenhang zum Employer Branding – die Auswirkungen der Personalmarketingstrategie auf das Recruiting.  & verändert: Passagen ausgelassen  \\
        ChatGPT 3 & Entwurf einer Gliederung für eine Hausarbeit zum Thema Recruiting  & unverändert \\
        Elicit 1 & Which elements should be included in an Employer Branding Plan?  & Passagen überarbeitet  \\
    \end{longtable}}
\else
    Es wurde keine KI verwendet.
\fi

% KI-Verzeichnis (verpflichtend)

% Rechtsquellenverzeichnis (sofern notwendig), z.B.:
%\addchap{Rechtsquellenverzeichnis}

% Anhangsverzeichnis und Anhang (sofern notwendig), z.B.:
%\addchap{Anhangsverzeichnis}
%...

% Vorlage für Praktikumsberichte
%
% Erklärung.tex
%
% Wortlaut gemäß Studien- und Prüfungsordnung/Anforderungsdokument 
% ("Selbstständigkeitserklärung") - nicht verändern.

\addchap{Selbstständigkeitserklärung}

Hiermit bestätige ich, dass ich die vorliegende Arbeit persönlich und selbständig verfasst und keine anderen als die angegebenen Quellen und Hilfsmittel verwendet habe. Alle Stellen, die wörtlich oder sinngemäß anderen Quellen entnommen wurden, sind als solche kenntlich gemacht. Die Zeichnungen, Abbildungen und Tabellen in dieser Arbeit sind von mir selbst erstellt oder wurden mit einem entsprechenden Quellennachweis versehen. Diese Arbeit wurde weder in gleicher noch in ähnlicher Form von mir an anderen Hochschulen zur Erlangung eines akademischen Abschlusses eingereicht.

\vspace{2cm}

% Ort und Datum sind in jedem Fall anzugeben (Werte aus Berichtinfos.tex)
\noindent
\begin{tabular}{@{}p{0.45\textwidth}@{\hspace{0.1\textwidth}}p{0.45\textwidth}@{}}
    \ort, \abgabedatum &
    \ifdefempty{\unterschrift}{}{\includegraphics[height=1.5cm]{\unterschrift}} \\
    \hrule & \hrule \\
    {\footnotesize Ort, Datum} & {\footnotesize \fullname}
\end{tabular}

\vspace{2cm}

\todo[inline]{Die Selbständigkeitserklärung ist mit der Unterschrift des Studierenden zu versehen (\texttt{\textbackslash unterschrift} in Berichtinfos.tex). Bei einem UPLOAD der Arbeit im HVS ersetzt der Authentifizierungsprozess die Unterschrift--diese ist dann optional.\\
In jedem Fall muss unter dem Text der ehrenwörtlichen Erklärung Ort und Datum angegeben werden.
}

% Selbstst�ndigkeitserkl�rung (verpflichtend)

\end{document}
%beendet das Dokument
